\cotefile{137}
\scanpage{73}
\uline{Ex 2} \quad Le ``complexe de De Rham \unsure{à|:} p. div''
est défini \deleted{comme} \unsure{canoniquement} (une \unsure{fois|fs}
connues les p.d.) comme \deleted{objet}
objet bigradué de $\deleted{\unsure{K}}\unsure{\mathcal R|\mathbb R|\mathcal Q}_{*}$, avec
structure d'algèbre bigraduée alternée
etc. Donc $F_C$ le transforme en
objet de la $=$ cat. On ne peut
affirmer encore (sans introduire
la ``structure p.d'' sur $F$) que
le complexe de De Rham associé à $F$
est alterné (seult antisym., i.e.,
\unsure{un signe|unité}). Le fait que
\deleted{\unsure{le} complexe de}
pour tt degré total donné ``$n$'', on trouve
une résolution de $h_*=\Psi_*^0$, flasque
(i.e. $\Sigma$-acyclique) sauf en degré
vert. maximum $n$ (où on trouve $\Psi_*^n
=\Lambda^n\Psi_*$) -- donc que c'est une résolution
$\Sigma$-acyclique tronquée, montre que
la \unsure{coho} de sa transformée par $F$
coïncide avec celle du $C^\bullet$ \deleted{\unsure{au fond}}\inserted{jusqu'en degré $n$}, et
\deleted{est nulle au dessus} -- mieux,
le complexe est homotope (avec
classe d'homotopie donnée) au
complexe tronqué \unsure{naïve|naïvem|naïf} de $C^\bullet$ --
en fait, on a un choix canonique.
(car $\overset{\inserted{\to\,\unsure{\deg\ \mathrm{tot}|\deg\leq n}}}{DR_*^{n\bullet}}\longrightarrow C^\bullet_*$ hom. \uline{unique}).
% archive sheet number: 72
\scanpage{74}
Il y a de $=$ un complexe de
De Rham-Sullivan \deleted{\unsure{en}}\inserted{si $k$ de} car. $0$.
% archive sheet number: 73
\scanpage{75}
\uline{Catégories} \uline{\unsure{hopolynômiales|polynômiales|polynomiales}}

Une \uline{catégorie polynômiale} est une catégorie $P$,
munie d'une sous-catégorie (pas pleine en général)
$P_1$, avec les axiomes suivants
\margin{Notation $P(X,Y)$ pour $\operatorname{Hom}_P(X,Y)$}

a) $P$ admet des produits finis, \deleted{et l'objet final est}

b) $\operatorname{Ob}P_1=\operatorname{Ob}P$, \deleted{les produits et pour deux objets
dans $P$, leur produit dans $P$ est un produit
dans $P_1$ (i.e. si $X,Y,Z\in\operatorname{Ob}P_1$, \deleted{$X,Y\in\operatorname{Ob}P$} \unsure{l'ensemble de}
\inserted{(i.e si $X,Y,Z\in\operatorname{Ob}P_1$, une} flèche $Z\to X\times Y$ de $P$ est dans $P_1$ ssi
$(\unsure{u,v|v,w}) : X\to Y\times Z$ est dans $P_1$ ssi} (donc $P_1$ est décrite
par un ens\unsure{.} $\mathrm{Fl}_1\subset\mathrm{Fl}(P)=\mathrm{Fl}$ contenant les identités,
et stable par composition.) On veut de plus que
\deleted{$\mathrm{Fl}_1$ contienne les isos, et que} pour $X,Y\in\operatorname{Ob}P$,
leur produit $X\times Y$ soit un produit ds $P_1$
(i.e. une flèche $u=(\unsure{u',u''|u,v|v,w}) : Z\to X\times Y$ est
``de degré $1$'' ssi ses composantes $X,Y$ le sont),
\inserted{enfin si $0$ est l'objet final, la fl $X\to0$ est de degré $1$ (donc $0$ objet final de $P_1$)}

c) $P_1$ est une catégorie additive

Cela implique que \deleted{pour tt objet l'on a
un hom de monoïdes,} les $\operatorname{Hom}_P(X,Y)$
\deleted{\[
\mathbb Z\longrightarrow\operatorname{End}(\mathrm{id}_{P_1})\longleftarrow\operatorname{End}(\mathrm{id}_P)
\]
\unsure{de plus} pr tt $\deleted{\unsure{A}}X,Y\in\operatorname{Ob}P$, \unsure{ou un} $X$}
\deleted{i.e pour tt $X$.}
sont des groupes abéliens (car les objets $\deleted{X}$\inserted{de $P$} sont
des objets groupes abéliens de $P$), dépendent
fonctorielt de $X$, mais pas de $Y$
(i.e. on a
\[
\boxed{(f+g)\circ h=fh+g\circ h}
\]
mais \deleted{\unsure{non en général}}\inserted{on \unsure{peut|va} avoir,} éventuellement
\[
f\circ(g+h)\ne f\circ g+f\circ h\quad !
\]

Il en résulte aussi que $\mathbb Z^\times$ opère sur $P$
i.e on a un homomorphisme multiplicatif
\[
\mathbb Z^\times\longrightarrow\operatorname{End}(P_1)\longleftarrow\operatorname{End}(P)
\]
qui à $n\in\mathbb Z$ et $X\in\operatorname{Ob}P_1=\operatorname{Ob}P$ associe $n\,\mathrm{id}_X$.
Ainsi $\mathbb Z$ opère dans $\deleted{\operatorname{Hom}_{\unsure{P}}}P(X,Y)$ par
$\mu_n(f)=f\circ(n\,\mathrm{id}_Y)$
% archive sheet number: 74
\scanpage{76}
d) Il existe une décomposition de $\deleted{\operatorname{Hom}}P(X,Y)$
en
\[
P(X,Y)=\bigoplus_{n\geq0}P_n(X,Y)
\]
telle que pour $f_n\in P_n(X,Y)$, et $\lambda\in\mathbb Z$, on ait
\[
f_n\circ(\lambda\mathrm{id}_X)=\lambda^nf_n \tag*{$*$}
\]
Plus généralement, si $X_1,\ldots,X_r,Y$ sont des
objets de $P$, on a une décomposition
\[
P(\deleted{\unsure{\prod_i X_i}}X_1\times\cdots\times X_r,Y)=
\bigoplus_{n=(n_1,\ldots,n_r)\in\mathbb N^r}P_n(X_1,\ldots,X_r;Y)\tag*{($*$)}
\]
avec, pour $f\in P_{n_1,\ldots,n_r}(X_1,\ldots,X_r;Y)$ et $\lambda=(\lambda_i)_{1\leq i\leq r}\in\mathbb Z^r$,
on ait
\[
f\circ(\lambda_1\mathrm{id}_{X_1},\ldots,\lambda_r\mathrm{id}_{X_r})
=\lambda_1^{n_1}\lambda_r^{n_r}f.
\]

% margin note written sideways
\margin{\uline{Ex} si $P$ est un groupe annulé par $p$, alors l'op. de $\mathbb Z$ par $\lambda\longmapsto\lambda\,\mathrm{id}_P$ et $\lambda\longmapsto\lambda^p\,\mathrm{id}_{\mathbb Z}$ est la $\underline{=}$, donc on peut concentrer au choix $P$ en \unsure{un|une} \unsure{en} degré $1$, ou $p$ !}

\uline{Remarques} Cette décomposition \inserted{n'}est \inserted{pas} \unsure{unique}, un
\deleted{\unsure{sorte de|retour de|vertu de}} canular ``arithmétique'' -- il le serait

\deleted{\uline{Lemme} Soit $P$ un \uline{groupe abélien}, et
\[
u:(\mathbb Z^\times)^r\longrightarrow\operatorname{End}(P)
\]
un hom de monoïdes.
\deleted{Pour tt $n=(n_1,\ldots,n_r)\in\mathbb Z$,}
Supposons qu'il existe une décomposition en
somme directe
\[
P=\sum_{n=(n_1,\ldots,n_r)\in\mathbb N^r} P_n\unsure{,}
\]
telle que pour $f\in P_n$\inserted{ et $\lambda=(\lambda_1,\ldots,\lambda_r)\in\mathbb Z^r$,} on ait
\[
\deleted{\unsure{\mu_1}}u(\lambda)\cdot f=\lambda_1^{n_1}\cdots\lambda_r^{n_r}f.
\]
Alors cette décomposition est unique, plus
précisément on a pour tt $n\in\mathbb N^r$
\[
P_n=\{f\in P\mid u(\lambda)f=\lambda_1^{n_1}\cdots\lambda_r^{n_r}f\quad\forall\lambda\in\mathbb Z^r\}
\]
Il suffit de voir que si $n'\ne n$ et si
$g_{n'}\in P_{n'}$ est tel que $\unsure{\mu|u}(\lambda)g_{n'}=\lambda_1^{n_1}\cdots\lambda_r^{n_r}g_{n'}\quad\forall\lambda\in\mathbb Z^r$,
alors $g_{n'}=0$. Or on aura aussi
\[
\deleted{\unsure{\lambda}}u(\lambda)g_{n'}=\lambda_1^{n'_1}\cdots\lambda_r^{n'_r}g_{n'}
\]}
% margin note written sideways; outer right edge cut by the scan
\margin{\unsure{Il} \unsure{si} \illegible{1}}
% archive sheet number: 75
\scanpage{77}
\deleted{donc
\[
(\lambda^n-\lambda^{n'})g_{n'}=0\qquad\forall\lambda\in\mathbb Z^r
\]
et il suffit de voir que les $\lambda^n-\lambda^{n'}$ engendrent
l'idéal $1$, ou mieux}

si on exigeait qu'on a, pour tt $\mathbb Z$-algèbre $k$,
des compositions entre les $P_k(X,Y)\simeq P(X,Y)\otimes_{\mathbb Z}k$,
et en exigeant la décomposition plus haut
valable \deleted{pour}\inserted{dans} tt $P_k(\prod X_i,Y)$, avec $\lambda\in(k^\times)^{\unsure{r|n}}$. On
va mettre comme \uline{structure supplémentaire}
la \uuline{donnée} de la décomposition $(*)$, avec

e) Condition \unsure{des degrés|de degré} pour la composition
\diagram{\begin{tikzcd}[ampersand replacement=\&,column sep=small]
X_1 \& {\times\cdots\times} \& X_r\arrow[r,"{(n_1,\ldots,n_r)}"] \& Y\\
{Z^1_1\times\cdots Z^1_{s_1}}\arrow[u,"{m^1_1,\ldots,m^1_{s_1}}"'] \& \& {Z^r_1,\ldots,Z^r_{s_r}}\arrow[u,"{m^r_1,\ldots,m^r_{s_r}}"] \&
\end{tikzcd}}
% slanted
\margin{Oubli : les hom de la sous-catégorie $P_1$ sont bien ceux de degré $1$ pour la décomposition $P(X,Y)=\bigoplus_i P_i(X,Y)$}
% slanted; first symbols on two lines touch the scan boundary
\margin{faut aussi si $X_r=0$ $\unsure{P|\mathcal P}_{n_1,\ldots,n_r}(X_1,X_2,X_r\deleted{\unsure{,\ldots}};M)=\unsure{0}$ si $n_r\ne0$, $P_{n_1,\ldots,n_{r-1}}(X_1,\ldots,X_{r-1})$ si $n_r=0$}
le composé est de degré $(n_i m^i_1,n_i m^i_2,\ldots,n_i m^i_{s_i})$
en $Z^i_1,\ldots,Z^i_{s_i}$ ; de $=$

f) condition \inserted{\unsure{symétrie|de symétrie}} \deleted{\unsure{pour}} la permutation des facteurs

g) $\forall X\in\operatorname{Ob}P$, l'application canonique déduite de
$X\longrightarrow0$
\[
P_0(0,Y)\longrightarrow P_0(X,Y)
\]
est bijective (les ``applications constantes'' \deleted{\unsure{canoniques}}\inserted{de $X$ dans $Y$}
``ne dépendent pas de $X$'')

Si $k$ est un anneau comm. et si $P_1$ est muni
d'une structure $k$-linéaire par
\[
k\unsure{-|\longrightarrow}\operatorname{End}(P_1)\longleftarrow\operatorname{End}(P)
\]
\deleted{on dit que}
telle que $(*)$ plus haut soit vrai pour
$\lambda\in k^r$ (non seulement $\lambda\in\mathbb Z^r$), on dit que
$P$ est $k$-polynomiale.
% archive sheet number: 76
\scanpage{78}
\uline{Relations avec op.} $\otimes$ \uline{et} $\Gamma^p$.

La donnée d'une cat. polynomiale \inserted{$P$} permet
de définir un foncteur triadditif
\[
P_1^\circ\times P_1^\circ\times P_1\longrightarrow\mathrm{Ab}
\]
par
\[
(X,Y;Z)\longmapsto P_{(1,1)}(X,Y;Z)
\]
(et de $=$ pour plus d'\uline{arguments} à gauche)

Si, pour $X,Y$ fixés, $Z$ variable, ce foncteur
est représentable, l'objet qui le représente est
noté $X\otimes Y$. Il dépend fonctorielt de
$X,Y$ ds $P_1$ (mais non pour $X,Y\deleted{\in\operatorname{Ob}P}$ considérés
comme dans $P$ !). \deleted{\unsure{On}} voit que $X\otimes Y$ existe
ssi $Y\otimes X$ existe, et ils sont can. isom. (\deleted{\unsure{fl|par}}
$s_{X,Y}$), on a $s_{X\otimes Y}\cdot s_{Y\otimes X}=\mathrm{id}_{Y\otimes X}$.

Soit $X,n\in\mathbb N$ fixés, si le foncteur
\[
Y\longmapsto P_n(X,Y)
\]
sur $P_1$ est représentbl, l'objet qui le représente
est noté $\Gamma^n X$ ; il \inserted{en} dépend aussi fonctorielt
de $X$ ds $P_1$. L'objet $\unsure{\Gamma_{\scriptscriptstyle\circ}^0|\Gamma^0}(X)$ existe ssi $\Gamma^0(0)$ existe
il est noté $\mathbb{1}$.

\unsure{Prouvons|Pouvons} que \deleted{\unsure{$\forall$}} si $\mathbb{1}$ existe, alors $\forall X\in\operatorname{Ob}\unsure{P|\Gamma|\mathcal P}$,
$\mathbb{1}\otimes X\simeq X$ i.e.
\[
P_{(1,1)}(\mathbb{1},X;Y)\underset{\inserted{\text{isom bifonctoriel}}}{\simeq}P_1(X,Y)(=\operatorname{Hom}_{P_1}(X,Y))\quad ?
\]
et plus particulièrement $\mathbb{1}\otimes\mathbb{1}\simeq\mathbb{1}$ i.e. $\forall Y$
\[
P_{1,1}(\mathbb{1},\mathbb{1};Y)\simeq P_1(\mathbb{1},Y)\quad ?
\]
% archive sheet number: 77
\scanpage{79}
On est obligé de renforcer la propriété universelle
de $\otimes$ et $\Gamma^p$, pour avoir les propriétés d'assoc. etc
qu'on veut.

Soit $X\times Y\longrightarrow\deleted{\unsure{P}}\deleted{\unsure{P}}T$ un morphisme de bidegré $(1,1)$,
alors pour tt $Z$ on a
\[
X\times Y\times Z\longrightarrow T\times Z
\]
d'où pour tt $M$
\[
P_{1,n}(T,Z;M)\longrightarrow P_{1,1,n}(X,Y,Z;M)
\]
On dit que $T$ est un produit \uline{tensoriel strict} de
$X$ et $Y$ si l'hom précédent est un iso,
\deleted{et \unsure{alors} $X\otimes Y=T$}
pour tt $n$ et tout $M$. Faisant $n=\unsure{1|0}$, \inserted{$Z=0$}, on trouve
\[
P_{1,\deleted{\unsure{1}}}(T,M)\simeq P_{1,1}(X,Y;M)
\]
donc $T$ est un produit tensoriel. C'est
en exigeant qu'on a des $\otimes$ \uline{stricts} \deleted{pour}\inserted{\unsure{de deux}} et
trois facteurs (\deleted{\unsure{qu'il}} il faut que la définition précédente
\inserted{au cas d'un produit \unsure{tensoriel de $n$ facteurs}}) que l'on
trouve \unsure{automatiqt} une structure
ass. comm. compatible avec
\[
\deleted{\operatorname{Hom}_{P_1}} P_1(X_1\otimes\cdots\otimes X_r;M)\simeq P_{1,\ldots,1}(X_1,\ldots,X_r;M)
\]
On sait qu'un objet unité (sous réserve d'une
donnée d'associativité) si elle existe, est unique.
\margin{?} Il faut examiner pourquoi ce devrait être $\Gamma^0(0)$.

Revenons aux $\Gamma^n$, soit $X=X_1\times\cdots\times X_r$,
tel que $\Gamma^n(X)$ existe. Considérons les
sous foncteurs $P_{\nu=(\nu_1,\ldots,\nu_r)}(X_1,\ldots,X_r;M)$
du foncteur \inserted{$P_1\longrightarrow(\mathrm{Ab})$} (cov.) représenté par $\Gamma^n(X)$,
supposons qu'ils soient aussi \unsure{corepr.|repr.} (automatique
si $P_1$ est karoubienne) soit $\Gamma^{\nu_1,\ldots,\nu_r}(X_1,\ldots,X_r)$
$=\Gamma^\nu(X_1,\ldots,X_r)$ l'\deleted{\unsure{foncteur}}\inserted{objet} qui le représente. $\deleted{\unsure{X_1}}$
On va définir \unsure{leur notion|leurs niveaux|sans réserve|sous réserve} d'existence des $\Gamma^{\nu_i}X_i$
\deleted{produit}\inserted{et des} tensoriel, \deleted{\unsure{en}} ou dans $\check{P}_1=\widehat{P_1^\circ}^{\,\circ}$ si.
% archive sheet number: 78
\scanpage{80}
on veut faire abstraction des questions de
représentabilité (sauf pour les $\Gamma^{\nu_i}(X_i)$)
\[
\Gamma^{\nu_1}(X_1)\otimes\cdots\otimes\Gamma^{\nu_r}(X_r)\longleftarrow\Gamma^\nu(X_1,\ldots,X_r)
\]
on aura un hom fonctoriel en
$M$ ds $P_1$ :
\[
P_{\nu}^{\unsure{P|\circ|\mathfrak{g}}}(X_1,\ldots,X_r;M)\longleftarrow P_{1,\ldots,1}(\Gamma^{\nu_1}X_1,\ldots,\Gamma^{\nu_r}X_r;M)
\]
par composition avec les morphismes canoniques
dans $P$
\[
X_1\xrightarrow{\gamma_{\nu_1}}\Gamma^{\nu_1}X_1,\ \ldots,\ X_r\xrightarrow{\gamma_{\nu_r}}\Gamma^{\nu_r}(X_r)
\]
et on suppose que c'est \inserted{tjrs} un iso.
Moyennant ceci, on a
\[
\Gamma^n(X_{\deleted{\unsure{i}}1}\times\cdots\times X_r)\xleftarrow{\sim}
\sum_{\substack{\nu=(\nu_i)\in\mathbb N^r\\\sum\nu_i=n}}
\underbrace{\Gamma^{\nu_1}(X_1)\otimes\cdots\otimes\Gamma^{\nu_r}(X_r)}_{\unsure{\Vert|=}\ \Gamma^{\nu_1,\ldots,\nu_r}(X_1,\ldots,X_r)}
\]

[ \uline{NB} Il suffirait de l'établir \unsure{avec les|sous les}
hypothèses relatives $r=2$ -- pour $r=2$ d'abord,
et le déduire ensuite pour $r$ \unsure{qcque|quelc.}.

En particulier, on trouve, \deleted{que} si on
considère le produit $X\times0$, pour $\mathbb{1}=\Gamma^0(0)$
\deleted{$\Gamma^1(X\times\unsure{1|0})\simeq$}
\[
X\otimes\mathbb{1}\simeq\Gamma^1(X)\otimes\Gamma^0(0)\simeq\Gamma^{1,0}(X,0)\simeq X
\]
\inserted{Car $P_{1,0}^{\unsure{*|P|\mathfrak{g}}}(X;0;M)\simeq P_1(X,M)$}
et de $=$ $\mathbb{1}\otimes X\simeq X$ \deleted{\unsure{$\mathbb{1}\otimes\mathbb{1}$ aussi}}
\deleted{\unsure{$X\otimes\mathbb{1}$}}

% slanted
\margin{\unsure{trouver} un \unsure{critère} simple \unsure{en genre}}
On suppose qu'il y a des produits tensoriels
stricts d'un \unsure{nb|nombre} fini de facteurs, et que
l'on a des $\Gamma^n X$ satisfaisant à la condition
plus haut.
% archive sheet number: 79
\scanpage{81}
Ainsi
\[
X\unsure{-|\longrightarrow|\longmapsto}\Gamma^*(X)=(\Gamma^i(X))_{i\geq0}
\]
est un foncteur (pas \inserted{\unsure{néct|ncés}} additif !)
\[
P_1\longrightarrow\operatorname{Grad}(P_1)
\]
qui est en fait un foncteur
``transformant sommes en $\otimes$''. Comme
tout objet de $P_1$ est un \inserted{objet} groupe comm.
pour l'opération $\oplus$, il s'ensuit que
dans chaque $\Gamma^*(X)$ il y a une structure
de $\otimes$-algèbre \inserted{graduée} comm., fonctorielle en $X$.
% slanted
\margin{Pour la $=$ raison $\Gamma^*X$ sera \unsure{est} aussi une cogèbre, \deleted{\unsure{donc}} et $=$ une \uline{bigèbre} ass. commut pour les $2$ structures.}

\deleted{En particulier}
On va définir \inserted{ds $P_1$} un hom, fonctoriel en $X$
\[
\Gamma^p(\Gamma^q(X))\longleftarrow\Gamma^{pq}(X)\tag*{($*$)}
\]
ou ce qui revient au $=$, un hom \deleted{\unsure{dans}}
de degré $pq$
\[
X\longrightarrow\Gamma^p(\Gamma^q(X))
\]
On prendra le composé
\diagram{\begin{tikzcd}[ampersand replacement=\&,column sep=large]
X\arrow[r,"{e_q^X}","{\substack{\text{hom uni}\\\text{de degré }q}}"'] \& \Gamma^q(X)\arrow[r,"{e_p^{\Gamma^q(X)}}","{\substack{\text{hom uni}\\\text{de degré }p.}}"'] \& \Gamma^p(\Gamma^q(X))
\end{tikzcd}}

On a la condition ``d'associativité''
\diagram{$\unsure{(*)|(\times)}\quad$\begin{tikzcd}[ampersand replacement=\&,column sep=large]
{\Gamma^{pqr}(X)}\arrow[r,dash,"{\alpha^X_{p,qr}}"]\arrow[d,"{\alpha^X_{pq,r}}"'] \& {\Gamma^p(\Gamma^{qr}(X))}\arrow[d,"{\Gamma^p(\alpha_{q,r}^{\unsure{\Gamma^r(X)|X}})}"]\\
{\Gamma^{pq}(\Gamma^r(X))}\arrow[r,"{\deleted{\unsure{\Gamma}}\alpha^{\Gamma^r(X)}_{p,q}}"] \& {\Gamma^p(\Gamma^q(\Gamma^r(X)))}
\arrow[from=1-1,to=2-2,phantom,"\text{comm.}" description]
\end{tikzcd}}
et d'\uline{unité} : l'isom \unsure{$\alpha$|.} est l'identité \unsure{si}
\unsure{l'on pose|l'on prend} $\Gamma^1 X=X$, $\Gamma^0 X=\mathbb{1}$
% archive sheet number: 80
\scanpage{82}
Inversement, soit $P_1$ \inserted{\unsure{une}} catégorie \inserted{$\otimes$}\unsure{-}additive \inserted{avec},
munie d'un \inserted{$\otimes$ \deleted{\unsure{ass}} \unsure{avec unité|avec objet unité|ACU strict}} foncteur de $(P_1,+)$
dans $(P_{1+}^{\mathrm{grad.}},\otimes)$, \deleted{\unsure{deux} \unsure{notif}}
\[
\Gamma^*:(P_1,\oplus)\longrightarrow(P_1^{\mathrm{grad}},\otimes)
\]
donc
\[
\Gamma^*(0)=\mathbb{1}\qquad\text{i.e.}\qquad
\Gamma^i(0)=\begin{cases}\mathbb{1}&\text{si }i=0\\0&\text{si }i>0\end{cases}
\]
\[
\Gamma^*(\deleted{M}X\oplus\deleted{N}Y)\simeq
\overbrace{\Gamma^*(X)\otimes\Gamma^*(\unsure{M|Y})}^{\text{avec degr \unsure{tot.|total}}}
\]
(isom \deleted{\unsure{de gr.|gr.}} donné, bifonctoriel en $X,Y$, compatible
avec contraintes d'unité, d'ass., de comm.)
avec
\[
\Gamma^i(M)=\begin{cases}0&\text{si }i<0\\\mathbb{1}&\text{si }i=0\\M&\text{si }i=1\end{cases}
\]
(isom de foncteurs en $M$), Supposons de
plus donné des hom de foncteurs
\[
\Gamma^{pq}\xrightarrow{\alpha_{pq}}\Gamma^p\deleted{\unsure{(}}\circ\Gamma^q
\]
(i.e.
\[
\Gamma^{pq}(X)\xrightarrow{\alpha_{pq}^X}\Gamma^p(\Gamma^q(X))
\]
)
satisfaisant la condition d'associativité
ci dessus, plus les conditions
évidentes pour \deleted{$p,q$} l'un des nbs $p,q$
égaux à $0$ ou $1$ :
\diagram{\begin{tikzcd}[ampersand replacement=\&,column sep=large]
{(**)} \& \Gamma^q\arrow[r,"{\alpha_{1,q}=\mathrm{id}}"]\arrow[dr,bend right=24,"{\alpha_{q,1}=\mathrm{id}}"'] \& {\Gamma^1\circ\Gamma^q=\mathrm{id}\circ\Gamma^q=\Gamma^q}\\
\& \& {\Gamma^q\circ\Gamma^1=\Gamma^1\circ\mathrm{id}=\Gamma^q}
\end{tikzcd}}
\diagram{\begin{tikzcd}[ampersand replacement=\&,column sep=large]
{(***)} \& \Gamma^0\arrow[r,"{\alpha_{0,q}=\mathrm{id}}"]\arrow[dr,bend right=24,"{\alpha_{q,0}\ \text{un }\text{\uuline{iso}}}"'] \& {\Gamma^0\circ\Gamma^q=\Gamma^0}\\
\& \& {\Gamma^q\circ\Gamma^0.}
\end{tikzcd}}
% archive sheet number: 81
\scanpage{83}
Les $\alpha_{p,q}$ permettent de définir des
\uline{\unsure{Définissons|Définitions}}.
\[
P_{\unsure{n|p}}(X,Y)=\operatorname{Hom}_{P_1}(\Gamma^n X,Y)
\]
\inserted{donc
\[
\left\{\begin{aligned}
P_1(X,Y)&=\operatorname{Hom}_{P_1}(X,Y)\\
P_0(X,Y)&=\operatorname{Hom}_{P_1}(\mathbb{1},Y)
\end{aligned}\right.
\]}
et des applications de composition
\[
\begin{aligned}
P_p(X,Y)\times P_q(Y,Z)&\longrightarrow P_{pq}(X,Z)\\
f\qquad g&\longmapsto g\circ f
\end{aligned}
\]
par
\diagram{\begin{tikzcd}[ampersand replacement=\&,column sep=small]
\Gamma^p X\arrow[r,"f"] \& Y \& \Gamma^qY\arrow[r,"g"] \& Z\\
{\Gamma^q(\Gamma^p(X))}\arrow[urr,"{\Gamma^q(f)}"]\arrow[r,dash,"\deleted{\unsure{\Gamma^p}}"'] \& \deleted{\unsure{\Gamma^q(Y)}}\arrow[rr,dash] \& \& T\\
{\Gamma^{pq}(X)}\arrow[u,"{\alpha_{q,p}^X}"] \& \& \& {g\circ f=g\Gamma^q(f)\alpha_{q,p}^X}
\end{tikzcd}}

Le diagramme comm. $(*)$ assure l'associativité
de la composition, $(**)$ \unsure{tout uniment} exprime que si \deleted{$g$ \unsure{resp.} $f$ est
de degré $1$,} \inserted{$p$ ou $q=1$,} alors le composé est celui
provenant de la bifonctorialité de
$P_p(X,Y)$ sur $P_1^\circ\times P_1$, enfin les compatibilités
$(***)$ signifient, pour la \unsure{deuxième|dernière|dernier}, que
\[
g\circ f=g\qquad\text{si $g$ de degré $0$}
\]
la deuxième signifie
\diagram{\begin{tikzcd}[ampersand replacement=\&]
|[draw,ellipse,inner sep=6pt]|{\overset{\deleted{\unsure{\Gamma}}}{\mathbb{1}}\xrightarrow{\substack{\alpha_{q,0}\\\sim}}\Gamma^q(\mathbb{1})}
\end{tikzcd}}
% slanted
\margin{\unsure{plus tard|pour plus tard}}
\unsure{on|je} l'avait oublié dans les axiomes
préliminaires -- il ne semble pas
\inserted{pas se déduire du \unsure{reste|restant}\unsure{.|\ldots}}
% archive sheet number: 82
\scanpage{84}
Enfin, si $X=\bigoplus_{1\leq i\leq r}X_i$, la décomposition
de $\Gamma^*(X)$ en \deleted{les} $\bigotimes\Gamma^*(X_i)$ donne
\deleted{une} décomposition des $P_n(\prod X_i,Y)$ en
des $P_{\nu_1\ldots\nu_r}(X_1,\ldots,X_r,Y)$.

On vérifie alors les axiomes
a) à g), et le fait qu'\deleted{il y a}
les produits tensoriels de $P_1$
sont \uline{stricts} \unsure{que|qq} $\unsure{\theta|0|\otimes}$ (nb quelcon. de
facteurs), que les $\Gamma^{\nu_i}$ existent et
satisfont la condition énoncée précédemment \dots.

\medskip\noindent\hfil\rule{4cm}{0.4pt}\hfil\par

Une catégorie \unsure{polynomiale|polynômiale} \inserted{$P$} est dite
\uline{stricte} si elle satisfait aux
conditions préliminaires permettant
de définir $\otimes$ et $\Gamma^*$ avec toutes les
propriétés envisagées. D'autre part,
une cat additive $P_{1}$ avec $\otimes$ et $\Gamma^*$
avec les structures décrites, est
dite une $\otimes$-pd-catégorie. Les
deux notions sont donc en \unsure{équivalence}.

On dit que $P$ est \deleted{\unsure{$k$-li}} une
catégorie polynômiale (resp $\otimes$-pd-cat)
$k$-linéaire ($k$ anneau comm)
si on a donné un hom d'anneau
\[
k\longrightarrow\operatorname{End}(\unsure{\mathrm{Id}|\mathrm{id}}_{P_1})\quad(\subset\mathrm{id}_P)
\]
% archive sheet number: 83
\scanpage{85}
tel que la décomposition de
$P_n(X_1\times\cdots\times X_r;Y)$ en des $P_{\nu_1\ldots\nu_r}(X_1,\ldots,X_r;Y)$
respecte les op. de $k$, i.e. $f\in P_{\nu_1\ldots\nu_r}$
implique $f\circ(\mathrm{id}_{\lambda_1}\times\cdots\times\mathrm{id}_{\lambda_r})=\lambda_1^{\nu_1}\cdots\lambda_r^{\nu_r}f$
pour tt $(\lambda_1,\ldots,\lambda_r)\in k^r$.

\medskip\noindent\hfil\shortstack[r]{\rule{3cm}{0.4pt}\\[-0.9ex]\rule{0.9cm}{0.4pt}}\hfil\par

\uline{Ex} La catégorie des $k$-modules
est une $\otimes$-pd-catégorie
De $=$ si $\unsure{\mathscr X|\mathscr E}$ un topos avec une
$k$-\deleted{\unsure{Alg.}}\unsure{Algèbre} $K$, la cat. des $K$-Modules
\deleted{et on peut} [ De $=$ on peut prendre des
$K$-Modules gradués, \deleted{\unsure{ent|par}} \inserted{par un groupe}
quelconque, \deleted{\unsure{laquelle} \unsure{soit}} on considère des
\margin{?}
produits $\Gamma^n(M^i)$ pour les degrés $ni$. ]

\medskip\noindent\hfil\shortstack[r]{\rule{3cm}{0.4pt}\\[-0.9ex]\rule{0.9cm}{0.4pt}}\hfil\par

Il y a diverses questions dans le cas
``abstrait''. Par exemple : quand
et comment \unsure{peuton|peut-on} \inserted{($\forall M\in\operatorname{Ob}P_1$)} sur
$\Gamma^*(M)$ une structure ``à p.d.'',
i.e. des hom de degré $1$
\[
\Gamma^p(\Gamma^q(M))\unsure{-|\longrightarrow}\Gamma^{pq}(M)
\]
satisfaisant aux axiomes des p.d. ?
% archive sheet number: 84
\scanpage{86}
Il y a une telle structure
canonique dans le cas de l'exemple.
Ou aussi : quand peut-on considérer
la composition des morphismes de $P$
\[
\deleted{\operatorname{Hom}_P}P(X,Y)\times P(Y,Z)-P(X,Z)
\]
($k$-linéaire en le deuxième argument)
comme une application linéaire
\[
P(Y,Z)\longrightarrow(\text{Appl. $k$-\unsure{polynomiales|polynomiale}})(P(X,Y)\to P(X,Z))
\]
(Ceci aussi est acquis dans l'exemple principal envisagé). La
construction doit se ramener pour
l'essentiel à \unsure{disposer|définir} de telles
structures suppl. sur les foncteurs $\Gamma^{\unsure{p|*|n}}$,
i.e. pour tt $X,Y$
\[
f\longmapsto\Gamma^p(f):\operatorname{Hom}_1(X,Y)\longrightarrow\operatorname{Hom}_1(\Gamma^p X,\Gamma^p Y)
\]
doit être défini comme ``application pol. de
degré $p$'' \unsure{entre} les $k$-modules
$\operatorname{Hom}_1(X,Y)$ et $\operatorname{Hom}_1(\Gamma^p X,\Gamma^p Y)$, et non
seulement comme application ensembliste
% archive sheet number: 85
